% ==================================================
% CHAPTER 8: DEPLOYMENT AND OPERATIONS
% ==================================================

\chapter{Deployment and Operations}
\label{ch:deployment}

\section{Infrastructure}

MemoryLane runs on a single AWS EC2 \texttt{c7i-flex.large} instance (2 vCPU,
4\,GB RAM plus 2\,GB swap) in the Mumbai region, with an Elastic IP and the
domain \texttt{memorylane.site}. Three Docker Compose services make up the
system (Figure~\ref{fig:deployment}).

\begin{figure}[H]
\centering
\begin{tikzpicture}[node distance=6mm and 10mm]
  \tikzset{vol/.style={db, minimum width=17mm, minimum height=8mm, font=\scriptsize, fill=white}}
  \node[comp, minimum width=31mm] (caddy) {\textbf{caddy}\\\scriptsize ports 80, 443\\\scriptsize certificates, web app,\\\scriptsize extension zip, \texttt{/api} proxy};
  \node[comp, right=9mm of caddy, minimum width=29mm] (backend) {\textbf{backend}\\\scriptsize FastAPI + bge-small\\\scriptsize no published port};
  \node[comp, right=9mm of backend, minimum width=28mm] (solr) {\textbf{solr} 9.7\\\scriptsize 1\,GB heap\\\scriptsize \texttt{127.0.0.1:8983}};
  \node[vol, below=8mm of caddy] (v1) {caddy\_data};
  \node[vol, below=8mm of backend, xshift=-9.5mm] (v2) {models};
  \node[vol, below=8mm of backend, xshift=9.5mm] (v3) {accounts};
  \node[vol, below=8mm of solr] (v4) {solr\_data};
  \draw[flow] (caddy) -- (backend);
  \draw[flow] (backend) -- (solr);
  \draw (caddy) -- (v1); \draw (backend) -- (v2); \draw (backend) -- (v3); \draw (solr) -- (v4);
  \begin{scope}[on background layer]
    \node[group, fill=gray!4, fit=(caddy)(solr)(v1)(v4), inner sep=9pt,
          label={[font=\small\bfseries]above:EC2 c7i-flex.large $\cdot$ Docker Compose}] (ec2) {};
  \end{scope}
  \node[box, left=12mm of caddy] (net) {Internet\\\scriptsize HTTPS};
  \draw[flow] (net) -- (caddy);
\end{tikzpicture}
\caption{Deployment: three containers and their volumes}
\label{fig:deployment}
\end{figure}

\begin{itemize}
  \item \textbf{Caddy} \cite{caddy} obtains and renews certificates
  automatically, redirects HTTP to HTTPS with HSTS, serves the built web app
  and the extension download, sets the Content Security Policy, and proxies
  \texttt{/api/*} to the backend.
  \item The Caddy image \emph{builds} the web app and zips the extension
  during \texttt{docker compose build}, so the files served always match the
  deployed code.
  \item Only ports 80 and 443 are open. The backend is reachable only
  through Caddy, and Solr listens on the server's localhost.
\end{itemize}

\section{Releasing a Change}

\begin{lstlisting}[language=bash]
cd ~/ucs503p-202627-MemoryLane && git pull origin master
cd code/src && docker compose up -d --build
docker compose logs -f backend
\end{lstlisting}

Schema changes are made first, by script, because Solr's automatic field
creation is off. Pages can be re-processed into a new core with the re-index
script, which never modifies the source core, and page vectors were
back-filled for older pages with a resumable script. Backups are EBS
snapshots of the instance's volume.

\section{Operations}

\begin{itemize}
  \item \textbf{Health}: \texttt{GET /api/health} reports Solr, the model and
  Gemini; Docker health checks restart unhealthy services, and every
  service restarts after a reboot.
  \item \textbf{Logs} are capped at three 10\,MB files per service and hold
  page ids, not URLs or queries.
  \item \textbf{Graceful degradation}: without Gemini, summaries are
  unavailable and everything else works; if email fails, reports wait in
  SQLite and are retried; if Cognito is unreachable, known sign-ins keep
  working.
  \item \textbf{Documentation} is published to GitHub Pages with MkDocs'
  \texttt{gh-deploy} command in strict mode, and the CI workflow checks the
  build on every change.
\end{itemize}

\section{Distribution of the Extension}

The extension is downloaded from the website as a zip built from the
deployed code and loaded unpacked in Chrome. Its version is bumped with
every change (2.4.2 at the time of writing), so users can see whether they
are up to date; the server's own checks protect against outdated
extensions in the meantime.
